% ============================================
% Johns Hopkins Letter Template
% Assistant Professor Cover Letter – Florida State University (Department of Economics)
% ============================================

\documentclass[11pt,letterpaper]{letter}

\usepackage{graphicx}
\usepackage{eso-pic}
\usepackage{geometry}
\usepackage{microtype}
\usepackage[hidelinks]{hyperref}

% ----- Page Setup -----
\geometry{left=25mm,right=25mm,top=25mm,bottom=25mm}

% ----- Logo Placement (foreground, first page only) -----
\newcommand{\PutJHULogoFG}[1]{%
  \AddToShipoutPictureFG*{%
    \AtPageUpperLeft{%
      \hspace{10mm}%
      \raisebox{-38mm}{%
        \includegraphics[width=6cm]{#1}%
      }%
    }%
  }%
}

% ----- Sender Information -----
\signature{Decory Edwards}
\address{Department of Economics \\ Johns Hopkins University \\ Baltimore, MD 21218 \\ \texttt{dedwar65@jh.edu}}

\begin{document}
\PutJHULogoFG{Images/jhulogo.png}

\begin{letter}{Search Committee \\
Department of Economics \\
College of Social Sciences and Public Policy \\
Florida State University \\
Bellamy Building, Room 288 \\
Tallahassee, FL 32306-2180, USA}

\opening{Dear Members of the Search Committee,}

I am writing to express my interest in the tenure-track Assistant Professor position in the Department of Economics at Florida State University, beginning in Fall 2026. I am a Ph.D.\ candidate in Economics at Johns Hopkins University, advised by Professors Christopher D.\ Carroll, Nicholas W.\ Papageorge, and Stelios Fourakis, and I expect to receive my degree in May 2026. My research fields are macroeconomics, applied microeconomics, and computational economics, with an emphasis on structural methods and empirical modeling.

My research agenda is unified by a focus on understanding the determinants of wealth inequality and the mechanisms through which heterogeneity in returns to wealth shapes aggregate outcomes. In my job market paper, I estimate the degree of heterogeneity in individual portfolio returns using micro data from the Survey of Consumer Finances and simulate a structural model in which households face idiosyncratic return risk. I show that allowing for persistent heterogeneity in returns substantially improves the model’s ability to match the empirical distribution of wealth, offering a tractable way to connect micro-level return dispersion to macroeconomic inequality.

In a second project, I use the Health and Retirement Study to study the relationship between interpersonal trust and portfolio performance. Building on the literature that finds a hump-shaped relationship between trust and income, I show that a similar relationship holds for individual returns to wealth measured in the HRS data, highlighting a behavioral channel through which social attitudes shape saving and investment outcomes. This work also provides new estimates of the persistent component of returns to net worth, which motivate the calibration of heterogeneity in my structural model.

The Department’s search for an applied structural economist resonates strongly with my research approach. My work combines theory-driven modeling with empirical estimation using micro and household-level data—techniques that complement the Department’s strengths in public, labor, health, innovation, industrial organization, and urban economics. I am particularly excited about contributing to the Department’s research on policy-relevant issues through the use of structural estimation and applied econometrics. My training in computational methods and experience working with large, heterogeneous data sets position me well to collaborate across these fields.

\textbf{Teaching.} My teaching experience centers on undergraduate macroeconomics and an upper-level macro policy seminar, where I have led weekly discussion sections and held regular office hours for classes ranging from 16 to 40 students. Consistent with my teaching statement, my approach is student-centered: I aim to help students ``convince themselves'' that they have moved from confusion to understanding by pairing clear scaffolding with structured office-hour coaching that builds trust and confidence. At Florida State, I am prepared to teach econometrics at both the undergraduate and graduate levels—including the first-year Ph.D.\ core econometrics sequence—as well as courses in macroeconomics, applied microeconomics, and related quantitative fields.

I believe I would be a strong fit because my computational and empirical toolkit allows me to (i) produce robust, data-backed structural and econometric estimates, (ii) build and validate models that speak to policy-relevant questions at the intersection of macro, public, and applied microeconomics, and (iii) contribute to a collaborative, research-active environment that values both methodological rigor and policy relevance. I am excited by FSU’s commitment to high-quality research, teaching excellence, and collegiality, and I would be honored to contribute to the Department’s mission.

Thank you for your consideration. I welcome the opportunity to discuss how my research and teaching could complement Florida State University’s Department of Economics.

\closing{Sincerely,}

\end{letter}
\end{document}
